\section{Matching the Current-Temperature Response of Differential GMOS}
As was previously discussed, the planned GMOS connection scheme which was
established so far consists of two 2T connected GMOS transistors, each connected
in series to a fixed value resistor. This connection method was evaluated and
it was established that a large offset readout in no gas measurements. This
hinders the ability of the readout circuit to extract the
relatively low signal fluctuations compared to the offset value. \\
In this chapter we will formalize the origins of the offset and explore ways
minimize it.

\subsubsection{Necessity to Minimize Offset Voltage}
The differential voltage signal resulting in the bridge topology can be fed
to the LIA sub-circuit in either DC or AC coupling. The signal at the input
of the LIA is sampled and then digitally processed to extract the required
LIA outputs. An important point to keep in mind is that the sampling circuit
has a finite and adjustable voltage range and a fixed ADC resolution. From
that fact, we can deduce that the minimal voltage step, and thus
sensitivity, is related to the selected voltage range. Another issue is
related to the usage of AC coupling. As stated by the datasheet, AC coupling
input filter applies a 1Hz HPF on the input signal. Since we are required to
sample very low frequency signals (in the $mHz$ range) we are therefore
required to use DC coupling sampling and minimize the offset effect as much
as possible to not archive clamping.

\subsection{Optimization of Current-Temperature Relation}
As per previous experiments, the GMOS transistors were connected in 2T
connection and in series to a resistors of similar value. This setup
introduces several contributing factors to the eventual DC offset at the
sampling input. For example, resistor mismatch causing the transistors to
operate on different operation points. The GMOS transistors themselves may
have slight process mismatch. \\
Each branch is a separate NMOS transistor connected in series to a resistor and fed by $V_{DD, i}(i=blind, active)$. Since the operation point of interest for gas sensing GMOS is the subthreshold, the drain current holds the following relation:

If we assume that each transistor is fed by an independent supply $V_{DD,i}$ we can write the value of the transistor current:

\begin{equation}
    I_{D,i}(T) = I_{0,i}\exp\!\left(\dfrac{q\,(V_{GS,i}-V_{T,i}(T))}{n_i\,k\,T}\right) 
\end{equation}

With the prefactor:

\begin{equation}
    I_{0,i} = \mu_{0,i} C_{ox,i} \dfrac{W_i}{L_i}\!\left(\dfrac{kT}{q}\right)^{\!2}\!(n_i-1)\Big(\dfrac{T_0}{T}\Big)^{\!2}
\end{equation}

The voltage threshold may drift with temperature change according to the relation:

\begin{equation}
    V_{T,i}(T)=V_{T,i}(T_0)+\alpha_i(T-T_0),\quad \alpha_i\equiv \dfrac{dV_T}{dT}\Big|_i
\end{equation}

The two current relations $I_{blind}(T)$ and $I_{active}(T)$ are exactly the
same if all Taylor coefficients match. However, since we can always assume
that for our case, we are operating on a small signal region, thus the only
relevant coefficients to match are the first and second order coefficients:

\begin{equation}
    I_{D,\mathrm{blind}}(T_0)=I_{D,\mathrm{active}}(T_0),\quad \mathrm{TCC}_{\mathrm{blind}}(T_0)=\mathrm{TCC}_\mathrm{active}(T_0)
\end{equation}

The temperature coefficient ($TCC$) at the operation point:

Writing down the KVL equation for each transistor at $\mathbf{T_0}$:

\begin{equation}
    V_{DD,i} = I_{D,i}(u_i)\,R_i + V_{T,i}(T_0) + u_i
\end{equation}

Therefore, the TCC becomes:

\begin{equation}
    \mathrm{TCC}_i \approx -\dfrac{1}{I_{D,i}}\dfrac{dI_{D,i}}{dT} = -\dfrac{q}{n_i k T}\!\left[\dfrac{dV_{T,i}}{dT}+\dfrac{V_{GS,i}-V_{T,i}}{T}\right] - \dfrac{2}{T}
\end{equation}

So, the \textbf{TCC is set entirely by the overdrive voltage} and device specific parameters. Forcing $\mathrm{TCC}_{\mathrm{active}}=\mathrm{TCC}_{\mathrm{blind}}$ uniquely pins down the relation:

\begin{equation}
    \dfrac{1}{n_1}\!\left[\alpha_1+\dfrac{u_1}{T_0}\right]=\dfrac{1}{n_2}\!\left[\alpha_2+\dfrac{u_2}{T_0}\right]
\end{equation}

Once both $u_i$ are fixed, both currents $I_{D,i}(u_i)$ are also fixed. It follows that \textbf{its impossible to force the same bias current} at the same time as the TCC in a 2T connection. The choise of resistors in that case plays no role since they can only adjust the voltage drop on the transistor.
